\section{群作用}


\textbf{本节用到的记号，先一次交代清楚}：

\begin{itemize}
\item $G$：群；$X$（或 $\Omega$）：被作用的集合，不要求是群。
\item $G\curvearrowright X$ 或 $\rho:G\times X\to X,\ (g,x)\mapsto g\cdot x$：$G$ 在 $X$ 上的作用，满足 $e\cdot x=x$ 与 $(gh)\cdot x=g\cdot(h\cdot x)$。
\item $\rho_g:X\to X,\ \rho_g(x)=g\cdot x$：由 $g$ 给出的置换；$\rho:G\to\operatorname{Sym}(X)$ 是群同态。
\item $\operatorname{Sym}(X)$：$X$ 上全体双射构成的对称群。
\item $\ker\rho=\{g\in G:g\cdot x=x\ \text{对所有}\ x\}$：作用的核；核为 $\{e\}$ 时称作用\textbf{忠实}。
\item $G\cdot x=\{g\cdot x:g\in G\}$：$x$ 的轨道；$G_x=\{g:g\cdot x=x\}$：$x$ 的稳定子（下一节展开）。
\item $Z(G)$：中心；$\cong$：同构；$\le$：子群；$\times$：直积。
\item $GL_n(F)$ 与 $A\cdot v=Av$：矩阵群作用在列向量上的标准例子。
\end{itemize}


\subsection{定义}


设 $G$ 是群，$X$ 是集合。一个左群作用是一个映射

\[
\rho:G\times X\longrightarrow X,
\qquad (g,x)\longmapsto g\cdot x,
\]

满足

\[
 e\cdot x=x,
\qquad
(gh)\cdot x=g\cdot(h\cdot x)
\]

对所有 $g,h\in G$、$x\in X$ 成立。第二条说明群乘法与“先后作用”一致。

对每个 $g\in G$，定义置换

\[
\rho_g:X\to X,
\qquad \rho_g(x)=g\cdot x.
\]

由于 $g^{-1}$ 给出逆映射，$\rho_g\in\operatorname{Sym}(X)$。于是每个作用等价于一个群同态

\[
\rho:G\longrightarrow\operatorname{Sym}(X),
\qquad g\longmapsto\rho_g.
\]

作用的核是

\[
\ker\rho=\{g\in G:g\cdot x=x\text{ 对所有 }x\in X\}.
\]

若核为平凡群，则称作用是\textbf{忠实的}。因此“群作用”可以理解为把抽象群表示成一群具体的置换。

\subsection{基本例子}


\subsubsection{对称群的自然作用}


$\operatorname{Sym}(X)$ 作用在 $X$ 上：对 $\sigma\in\operatorname{Sym}(X)$ 和 $x\in X$，令 $\sigma\cdot x=\sigma(x)$。

\subsubsection{左正则作用}


群 $G$ 通过左乘作用在自身上：

\[
 g\cdot x=gx.
\]

该作用是忠实的，因为若 $g$ 固定所有 $x$，特别固定单位元 $e$，就有 $g=ge=e$。因此得到 Cayley 嵌入

\[
G\hookrightarrow\operatorname{Sym}(G).
\]

这说明每个群都同构于某个对称群的子群。

\subsubsection{共轭作用}


$G$ 作用在自身上，定义为

\[
 g\cdot x=gxg^{-1}.
\]

这个作用对应的同态就是 $g\mapsto c_g$，其核为 $Z(G)$。

\subsubsection{线性作用}


若 $G\le GL_n(F)$，则 $G$ 作用在向量空间 $F^n$ 上：

\[
 A\cdot v=Av.
\]

这是“群作用”把代数结构与几何对象连接起来的典型例子。

\subsection{作用的图像}


对固定 $x\in X$，作用给出一个轨道映射

\[
\alpha_x:G\longrightarrow X,
\qquad g\longmapsto g\cdot x.
\]

它的像是 $x$ 能够到达的所有位置；它的纤维将由稳定子陪集描述。这个观察是轨道--稳定子定理的核心。

\subsection{记忆点}


\begin{itemize}
\item 群作用的两条公理分别对应单位元和乘法相容。
\item 作用就是一个到对称群的同态；忠实性就是核平凡。
\item 左正则作用给出 Cayley 定理；共轭作用给出 $G/Z(G)\cong\operatorname{Inn}(G)$。
\end{itemize}


\subsection{参考对照}


\begin{itemize}
\item [DF] §1.7、§2.2：Group Actions、Stabilizers and Kernels。
\item [DS] §14.1.1–14.1.4：群作用、作用同态与 Cayley 定理。
\end{itemize}
